\section{Introduction}

Developers are increasingly using Large Language Models to write and review code~\citep{lee2025hallucination,MAHADEHASAN2026112815}. Most of these models are so large that they can only run on external computers~\citep{nguyen2026noirprivacypreservinggenerationcode}. This raises concerns about usage limits and cost, as well as keeping important information safe~\citep{nguyen2026noirprivacypreservinggenerationcode,shen2025gpiot}. On the other hand, Small Language Models are cheaper to run and can work on a local computer~\citep{MAHADEHASAN2026112815,shen2025gpiot}. This helps reduce security risks and reduces dependence on external services~\citep{nguyen2026noirprivacypreservinggenerationcode,shen2025gpiot}. These small models aren't as good as the big ones at writing code~\citep{MAHADEHASAN2026112815}. So the code they create might look plausible. Behaves incorrect. This is called Code Hallucination~\citep{pavel2025hallucination,lee2025hallucination,tian2025codehalu}. It's also hard to find these mistakes~\citep{andriushchenko-etal-2026-efficient,lee2025hallucination} because the models don't understand how code works the way a real programmer does. They just learn how to write code by looking at lots of examples. So writing code isn't the same as knowing all the details and the problems it might have.

In this work, we use the term \emph{Code Hallucination} to describe the phenomenon in which the model generates plausible but incorrect code, whose flaws translate into observable errors at an execution level~\citep{pavel2025hallucination,lee2025hallucination,tian2025codehalu}. Following previous work on errors and hallucinations in code generation~\citep{lee2025hallucination,dou2025whatswrong,tian2025codehalu}, we define 3 error levels, each associated with an observable execution level.\emph{Syntax Error} refers to hallucinations about the grammatical structure of the language that result in syntax errors that prevent compilation~\citep{lee2025hallucination,dou2025whatswrong}.\emph{Runtime Error} is a hallucination in which the model makes mistakes in execution logic and specific code knowledge, leading to failures during execution~\citep{lee2025hallucination,dou2025whatswrong}.\emph{Functional Error} is a hallucination that deviates from functional requirements and can be detected through manual checks or automated tests~\citep{lee2025hallucination,dou2025whatswrong,tian2025codehalu}, and it can be dangerous because it can be subtle enough to go unnoticed in poorly controlled environments.

Our central idea, therefore, is to study SLMs not only as code generators, but also to evaluate their real ability to analyze code correctness. We investigated whether various SLMs can recognize whether code is well-written or contains errors, and the degree of those errors, including in code the model generates itself. We also investigated how different forms of targeted training can improve these models' ability for this task. We compared these models with static tools to understand how they can serve as a useful alternative in real software environments, identify flaws without execution, and how combining static tools with the models can improve both their ability to analyze code and to correct it. Finally, we compared SLMs with larger LLMs to understand the difference between models that can run locally on a home computer and those that run on supercomputers.

To conduct this investigation, we used the EvalPlus benchmark, which provides some basic Python programming tasks and runs on various LLMs, primarily SLMs. We analyzed their code capabilities and, based on the execution of the generated code, measured how often models generated code that passed all available test cases and the level of execution achieved. After collecting hallucination data, we used the models to generate evidence-based explanations of generation errors for model training, using a variation of target parameters and metrics. Finally, py\_compile and pyright were used for a static analysis of syntax and execution errors.

The results highlight a divergence between error generation and error-detection capabilities. Most tested SLMs achieved 70-80\% accuracy in code generation, surpassing the larger models used in this study. However, the main problem with all tested models in code generation is incomplete task understanding, which results in code gaps and missing logic needed to fully solve the problem. In error detection, no model was able to accurately detect all classes, especially the \emph{Runtime Error} class, in which the model with the highest recall rate was 32\%, and the median among the models was only 14.6\%, highlighting the models' shortcomings in understanding code execution in practice. However, even with this limitation, the models still detect runtime failures better than static tools, showing they can detect runtime failures without execution and that combining them with static tools for manual compilation verification can make detection more robust. Finally, we found that specialization can substantially improve detection of these errors; however, the models may develop a pronounced bias due to class imbalance.

The main contributions of this work are:

\begin{itemize}

    \item An empirical evaluation of SLMs for both code generation and error detection under a common execution-based reference, distinguishing Syntax, Runtime, and Functional Errors through multilabel classification.

    \item An analysis of \emph{self-error detection}, investigating whether models can recognize errors in code they generated themselves and examining the relationship between generation and evaluation capabilities.

    \item A specialized SLM-based detector trained with execution-derived labels and teacher-generated explanations, together with an analysis of how different optimization objectives affect the detector's behavior across error categories.

    \item An investigation of the relationship between model-based reasoning and deterministic static analysis, evaluating static techniques independently, as supporting evidence for error detection, and as feedback during iterative code generation.

    \item A comparison between SLMs and larger models, assessing whether a clear gap exists in error generation, error detection, and tool assistance, to determine whether SLMs can be a viable alternative for these tasks.

\end{itemize}

This work is divided into sections that initially review how the literature has addressed code hallucinations, error detection, and specialization methods. Next, it describes the methodology and experimental setup. The results section presents the main findings on the SLM's code generation and evaluation ability; error-detection specialization; comparison and use of static tools; and, finally, a comparison of these abilities with larger models. The article concludes by discussing factors that may affect the results, summarizing the main findings, and suggesting areas for future research.
